\documentclass{beamer}

% You can also use a 16:9 aspect ratio:
%\documentclass[aspectratio=169]{beamer}
\usetheme{TACC16}
\usepackage{graphicx}

% It's possible to move the footer to the right:
%\usetheme[rightfooter]{TACC16}


% Detailed knowledge of application workload characteristics can
% optimize performance of current and future systems. This may sound
% daunting, with many HPC data centers hosting over 2,000 users running
% thousands of applications and millions of jobs per month.  XALT is an
% open source tool developed at the Texas Advanced Computing Center
% (TACC) that collects system usage information to quantitatively report
% how users are using your system. This session will explore the
% benefits of detailed application workload profiling and how the XALT
% tool has helped leading supercomputing sites unlock the power of their
% application usage data.

%% page 
%\begin{frame}{}
%  \begin{itemize}
%    \item
%  \end{itemize}
%\end{frame}
%
%% page 
%\begin{frame}[fragile]
%    \frametitle{}
% {\tiny
%    \begin{semiverbatim}
%    \end{semiverbatim}
%}
%  \begin{itemize}
%    \item
%  \end{itemize}
%
%\end{frame}



\begin{document}
\title[XALT]{Protecting XALT from Users}
\author{Robert McLay}
\date{March 17, 2022}

% page 1
\frame{\titlepage}

% page 2
\begin{frame}{XALT: Outline}
  \center{\includegraphics[width=.9\textwidth]{XALT_Sticker.png}}
  \begin{itemize}
    \item XALT is linking with every program that runs on the system
    \item Users will occasionally make mistakes
    \item Need to protect XALT from user mistakes
    \item Show three protection examples 
  \end{itemize}
\end{frame}

% page 3
\begin{frame}{Three examples of protection}
  \begin{itemize}
    \item User's bug hidden by zeroed memory initially
    \item User's mixing Fortran routine with C library routines badly
    \item XALT expecting well managed memory heap.
  \end{itemize}
\end{frame}


% page 4
\begin{frame}[fragile]
    \frametitle{How XALT works}
 {\tiny
    \begin{semiverbatim}
#include <stdio.h>
void myinit(int argc, char **argv)
\{ printf("This is run before main()\textbackslash{}n"); \}
void myfini()
\{ printf("This is run after main()\textbackslash{}n"); \}

__attribute__((section(".init_array"))) __typeof__(myinit) *__init = myinit;
__attribute__((section(".fini_array"))) __typeof__(myfini) *__fini = myfini;
    \end{semiverbatim}
}
  \begin{itemize}
    \item my\_docs/22/xalt\_monthly\_mtg\_2022\_03\_17/code/bad\_memory/ex1
  \end{itemize}
\end{frame}

% page 5
\begin{frame}[fragile]
    \frametitle{How XALT works (II)}
 {\small
    \begin{semiverbatim}
\% cat try.c

#include <stdio.h>
int main()
\{
  printf("Hello World!\textbackslash{}n");
  return 0;
\}

    \end{semiverbatim}
}
\end{frame}


% page 6
\begin{frame}[fragile]
    \frametitle{How XALT works (III)}
 {\small
    \begin{semiverbatim}
\$ ./try

Hello World!

\$ LD\_PRELOAD=./libxalt.so  ./try
This is run before main()
Hello World!
This is run after main()
    \end{semiverbatim}
}
  \begin{itemize}
    \item my\_docs/22/xalt\_monthly\_mtg\_2022\_03\_17/code/bad\_memory/ex1
  \end{itemize}

\end{frame}

% page 7
\begin{frame}{How XALT works (IV)}
  \begin{itemize}
    \item The xalt runtime library attaches to every program run on
      your system!
    \item I sometimes feel like I'm a developer on every program team.
    \item XALT programs are a in the same namespace as the user's
      program (UGH!)
  \end{itemize}
\end{frame}



% page 8
\begin{frame}[fragile]
    \frametitle{Hiding XALT routine names from users}
 {\tiny
    \begin{semiverbatim}
\% nm \$LD\_PRELOAD| grep \_\_XALT\_build
0000000000009e80 T \_\_XALT\_buildEnvT\_xalt\_1\_5
000000000000a840 T \_\_XALT\_buildUserT\_xalt\_1\_5
00000000000163a0 T \_\_XALT\_buildXALTRecordT\_xalt\_1\_5
    \end{semiverbatim}
}
  \begin{itemize}
    \item XALT routine names are hidden by macros supplied in xalt\_obfuscate.h
  \end{itemize}

\end{frame}


% page 9
\begin{frame}{User's bug hidden by initially zeroed memory}
  \begin{itemize}
    \item Initially all memory is zeroed before program starts
    \item Note that pointer zero, integer zero and float zero are all
      zero bits
    \item Link lists require a NULL pointer at end of list.
    \item Used memory is {\color{red} \emph{NOT}} zeroed for you in C.
    \item User's program work w/o XALT, Failed with XALT.
  \end{itemize}
\end{frame}

% page 10
\begin{frame}[fragile]
    \frametitle{Example code clean/used memory}
 {\tiny
    \begin{semiverbatim}
\% cat try.c

#include <stdio.h>
#include <stdlib.h>
#define SZ 1000
int main()
\{
  int *a = (int *) malloc(SZ*sizeof(int));
  printf("Hello World! a:%d\textbackslash{}n",a[0]);
  return 0;
\}

\% cat xalt.c

#include <stdio.h>
#include <stdlib.h>
#include <string.h>
#define SZ 1000
void myinit(int argc, char **argv)
\{
  int i;
  int *a = (int*) malloc(SZ*sizeof(int));
  for (i = 0; i < SZ; ++i) a[i] = 15; 
  free(a);
  printf("This is run before main()\textbackslash{}n");
\}
\_\_attribute\_\_((section(".init\_array"))) \_\_typeof\_\_(myinit) *\_\_init = myinit;
    \end{semiverbatim}
}

  \begin{itemize}
    \item my\_docs/22/xalt\_monthly\_mtg\_2022\_03\_17/code/bad\_memory/ex2
  \end{itemize}
\end{frame}

% page 11
\begin{frame}[fragile]
    \frametitle{Example code clean/used memory(II)}
 {\small
   \begin{semiverbatim}
\% ./try

Hello World! {\color{blue} a:0}

\% LD\_PRELOAD=./libxalt.so  ./try  ; echo
This is run before main()
Hello World! {\color{red} a:15}
This is run after main()
    \end{semiverbatim}
}
  \begin{itemize}
    \item my\_docs/22/xalt\_monthly\_mtg\_2022\_03\_17/code/bad\_memory/ex2
  \end{itemize}
\end{frame}

% page 12
\begin{frame}{XALT Fix: zero memory before free()}
  \begin{itemize}
    \item To protect XALT from broken user code
    \item XALT in myinit() zero's memory before free
    \item Note that non-MPI tracking does little allocation
    \item MPI tasks $>$ 127 init record $\Rightarrow$ much allocation
  \end{itemize}
\end{frame}



% page 13
\begin{frame}[fragile]
    \frametitle{XALT Fix: zero memory before free()}
 {\tiny
    \begin{semiverbatim}
\% cat try.c

#include <stdio.h>
#include <stdlib.h>
#define SZ 1000
int main()
\{
  int *a = (int *) malloc(SZ*sizeof(int));
  printf("Hello World! a:%d\textbackslash{}n",a[0]);
  return 0;
\}

\% cat xalt.c

#include <stdio.h>
#include <stdlib.h>
#include <string.h>
#define SZ 1000
void myinit(int argc, char **argv)
\{
  int i;
  int *a = (int*) malloc(SZ*sizeof(int));
  for (i = 0; i < SZ; ++i) a[i] = 15; 
  {\color{blue}{}memset((void *) a, 0, SZ*sizeof(int))};
  free(a);
  printf("This is run before main()\textbackslash{}n");
\}
\_\_attribute\_\_((section(".init\_array"))) \_\_typeof\_\_(myinit) *\_\_init = myinit;
    \end{semiverbatim}
}

  \begin{itemize}
    \item my\_docs/22/xalt\_monthly\_mtg\_2022\_03\_17/code/bad\_memory/ex3
  \end{itemize}
\end{frame}

% page 14
\begin{frame}[fragile]
    \frametitle{XALT Fix: zero memory before free() (II)}
 {\small
   \begin{semiverbatim}
\% ./try

Hello World! {\color{blue} a:0}

\% LD_PRELOAD=./libxalt.so  ./try  ; echo
This is run before main()
Hello World! {\color{blue} a:0}
This is run after main()
    \end{semiverbatim}
}
  \begin{itemize}
    \item my\_docs/22/xalt\_monthly\_mtg\_2022\_03\_17/code/bad\_memory/ex3
  \end{itemize}
\end{frame}


% page 15
\begin{frame}[fragile]
    \frametitle{Protecting XALT from Fortran mixed with C programs badly}
 {\tiny
    \begin{semiverbatim}
\% cat msg.f90
subroutine msg
   print *, "Hello World!"
end subroutine msg

\% nm try | grep msg
00000000000011c7 T msg_
    \end{semiverbatim}
}
  \begin{itemize}
    \item Normally Fortran routines get a trailing underscore
      when compiled
    \item This can be disabled:
    \item Fortran: -fno-underscoring
    \item ifort:    -assume nounderscore
    \item Can make mixing C/Fortran easier
    \item Also make collisions with C library easier
  \end{itemize}
\end{frame}

% page 16
\begin{frame}{XALT uses libuuid}
  \begin{itemize}
    \item libuuid.so is used to get a unique identifier
    \item It uses libc's random()
    \item Can't have two routines named random() 
    \item my\_docs/22/xalt\_monthly\_mtg\_2022\_03\_17/code/random/ex3


  \end{itemize}
\end{frame}

% page 17
\begin{frame}[fragile]
    \frametitle{Collision over random() routine}
 {\tiny
    \begin{semiverbatim}
\% cat try.f90

program tryMe
   implicit none 
   real*8 d, random
   print *, "Hello World!"
   d = {\color{blue}{}random(1.0, 2.0. 3.0)}
   print *, "d: ",d
end program tryMe

\% cat random.f90

real*8 function {\color{blue}{}random(a, b, c)}
   implicit none
   real*8 a, b, c
   print *, "In random(a, b, c)"
   random = a*b + c
end function random

\% cat xalt.c
#include <stdio.h>
#include <stdlib.h>
void myinit(int argc, char **argv)
\{
  long int a;
  printf("This is run before main()\textbackslash{}n");
  a = {\color{red}{}random()};
  printf("called random(): a: %ld\textbackslash{}n",a);
\}
__attribute__((section(".init_array"))) __typeof__(myinit) *__init = myinit;

\end{semiverbatim}
}
\end{frame}

% page 18
\begin{frame}[fragile]
    \frametitle{Collision over random() routine (II)}
 {\tiny
    \begin{semiverbatim}
\% ./try

 Hello World!
 In random(a, b, c)
 d:    5.0000000000000000     

\% LD_PRELOAD=./libxalt.so  ./try  ; echo
This is run before main()
 In random(a, b, c)
Segmentation fault
    \end{semiverbatim}
}
  \begin{itemize}
    \item The linker chooses the user's Fortran random() instead of the C lib
      random()
    \item The segfault happens because the Fortran random() expects 3
      arguments
    \item the random() call in xalt.c passes none.
  \end{itemize}

\end{frame}

% page 19
\begin{frame}{How to fix this issue}
  \begin{itemize}
    \item Other Fortran program might do the same thing
    \item Trick: Use dlopen()/dlsym() to dynamically link in libuuid.so
    \item At this point libuuid.so can't ``see'' the Fortran random()
      routine
    \item This trick solves many problems with libuuid
  \end{itemize}
\end{frame}

% page 20
\begin{frame}{XALT is still susceptible to similar issues}
  \begin{itemize}
    \item XALT is now protected from a user's random() function
    \item But XALT is vulnerable some Fortran code replacing a c
      library routine
    \item We will just have to fix them as they come up
  \end{itemize}
\end{frame}


% page 21
\begin{frame}[fragile]
    \frametitle{Protecting XALT from badly managed memory heap}
 {\small
    \begin{semiverbatim}
void my\_free(void *ptr,int sz)
\{
  if (s\_start\_record \&\& ptr != NULL)
    \{
      memset(ptr, '\textbackslash{}0', sz);
      free(ptr);
    \}
\}
    \end{semiverbatim}
}
  \begin{itemize}
    \item Reporting an end record in myfini() requires memory
      allocations
    \item However some user programs can leave the heap broken
    \item XALT replaces free() with my\_free()
    \item Memory is only freed for a start record.
  \end{itemize}
\end{frame}

% page 22
\begin{frame}{Conclusions}
  \begin{itemize}
    \item XALT has matured greatly from working with user programs
    \item Since the XALT library is in the same namespace as the user code 
    \item There is always a risk of routine collision.
  \end{itemize}
\end{frame}


% page 23
\begin{frame}{Future Topics?}
  \begin{itemize}
    \item Other protection of XALT from users
    \item Recent changes to importing json records
    \item Others?
  \end{itemize}
\end{frame}
%

%
\begin{frame}{License}
\footnotesize
\textcopyright The University of Texas at Austin, \the\year\\
\vspace{0.5 cm}
This work is licensed under the Creative Commons Attribution Non-Commercial 3.0 Unported License. To view a copy of this license, 
visit http://creativecommons.org/licenses/by-nc/3.0/\\
\vspace{0.5 cm}
When attributing this work, please use the following text: 
"Designing and Administering Large-scale Systems", Texas Advanced Computing Center, \the\year. Available under a Creative Commons Attribution Non-Commercial 3.0 Unported License.
\newline
\newline
\includegraphics[width=1cm]{./figures/cc.png}
\end{frame}

\end{document}
